\chapter{Evaluarea implementării}
\label{chapter:evaluare}

\section{Introducere}
Obiectivele stabilite pentru platformă au fost atinse, iar evaluarea implementării arată un succes în ceea ce privește funcționalitatea, 
performanța și experiența utilizatorului.

\section{Evaluarea back-endului}
\label{subsec:evaluare-backend}

\subsection{Testarea back-endului}
\label{subsubsec:evaluare-backend-testare}

Testarea back-endului s-a concentrat pe validarea funcționalităților API-urilor dezvoltate pentru interacțiunea cu baza de date
și manipularea datelor. Pentru aceasta, au fost utilizate soluții precum Postman și inspecția manuală a codului. În plus, au fost
testate diverse scenarii de utilizare pentru a verifica integritatea fluxului de date între componentele back-endului. Acestea includ
inclusiv inserarea, modificarea și ștergerea de date în MongoDB și performanța acestora.

De exemplu s-a testat API-ul pentru accesarea și manipularea formulelor a fost testat pentru răspunsuri rapide de sub
200ms pentru 1000 de cereri simultane.

\subsection{Monitorizarea back-endului}
\label{subsubsec:evaluare-backend-monitorizare}

Pentru monitorizarea back-endului, s-au utilizat instrumente de monitorizare și colectare precum Prometheus și Grafana pentru a
urmări performanța sistemului și a verifica orice problemă de scalabilitate sau fiabilitate. Datele colectate au arătat o utilizare 
constantă de resurse, cu peste 95\% din cereri servite în mod eficient, fără a depăși limitele impuse de infrastructura curentă.

\subsection{Evaluarea performanței back-endului}
\label{subsubsec:evaluare-backend-performanta}

Performanța back-endului a fost evaluată pe baza unui set de scenarii de utilizare. Scenariile au inclus de exemplu 100 de cereri simultane
la un API de calcul al formulelor pentru volum.

Testele au arătat o latență medie de răspuns de 250ms pe cereri simple și sub 2 secunde pentru încărcarea fișierelor mari. În ceea ce
privește scalabilitatea, s-au făcut teste de încărcare pentru a evalua comportamentul sistemului la mai mult de 1000 de cereri simultane,
iar infrastructura a rămas stabilă, fără scăderi semnificative ale performanței chiar și testat în mediu de dezvoltare cu resurse locale limitate.

Alt nivel de testare în timpul dezvoltării a fost prin \textbf{teste unitare},
rulate local cu \verb|go test|. Folosind pachetul \texttt{httptest} am
simulat cereri HTTP către ruterul unui serviciu din interior şi am putut
itera rapid și prinde bug-urile, fără să reconstruim imaginea Docker
sau să dăm deploy în Kubernetes de fiecare dată.

Rularea comenzii \verb|go test ./...| rulează orice test rapid,
oferind feedback aproape instant asupra
regresiilor de exemplu din transormatorul de cod, parserul sau din logica de
gestionare și propagare a erorilor de compilare către frontend.

\section{Evaluarea front-endului}
\label{subsec:evaluare-frontend}

\subsection{Testarea front-endului}
\label{subsubsec:evaluare-frontend-testare}

Testarea front-endului a fost realizată prin testare manuală a interacțiunilor cu utilizatorul pe diverse dispozitive. 
S-au testat funcționalitățile de navigare, interactivitatea componentelor (de exemplu, meniurile rotative, butoanele de creare a quiz-urilor),
și integrarea corectă a 3D-ului în scene. Testele au arătat că 95\% dintre interacțiuni au fost realizate fără erori, iar platforma a funcționat
decent pe dispozitivele de testare.


\subsection{Monitorizarea front-endului}
\label{subsubsec:evaluare-frontend-monitorizare}

Majoritatea paginilor se încarcă complet în mai puțin de 2 secunde pe majoritatea dispozitivelor, iar utilizarea memoriei este sub 
150MB pentru majoritatea scenelor 3D care încarcă obiecte și texturi.

\subsection{Evaluarea performanței front-endului}
\label{subsubsec:evaluare-frontend-performanta}

Performanța front-endului a fost evaluată prin testarea mai multor scenarii de utilizare, inclusiv interacțiuni cu animații 3D complexe 
și manipularea lor. Testele au arătat că platforma susține animații 3D simultane cu un cadru constant de 60fps, iar timpii de 
încărcare au fost sub 1,5 secunde pentru majoritatea paginilor.

\section{Testarea infrastructurii / platformei}
\label{subsec:evaluare-infrastructura}

S-au utilizat teste de scalabilitate și failover, iar platforma a demonstrat un comportament robust și o scalabilitate bună
în fața unui număr mare de cereri. De exemplu, un scenariu de testare a implicat 1000 de cereri simultane, iar sistemul a rămas stabil,
fără erori sau timp de nefuncționare. În plus, în caz că apar probleme de performanță, se pot scala serviciile în mod dinamic, adăugând
replici direct din scalarea oferită de Kubernetes pentru a răspunde la cerințele crescute.


\section{Probleme întâmpinate}
\label{subsec:probleme}

\subsection{Backend \& Kubernetes Ingress}
Una dintre principalele provocări întâmpinate a fost accesibilitatea Kong Ingress Controller în mediul local, care nu este disponibilă implicit.
Soluția implementată a fost folosirea unui port-forward pentru a expune Ingress-ul. Totodată, expunerea microserviciilor prin Kong a 
necesitat sincronizarea corectă a valorilor pentru containerPort și targetPort, configurarea serviciilor de tip ClusterIP și adnotările 
specifice pentru Ingress. Ajustările făcute au asigurat că traficul extern poate accesa microserviciile într-un mod stabil și controlat.

\subsection{Accesibilitate CI/CD Build \& Rollout Workflow}
Fluxul de lucru pentru CI/CD a fost un alt punct dificil de gestionat. Fiecare schimbare în cod presupunea reconstruirea și încărcarea
imaginii containerului, urmată de un proces de rollout. Pentru a automatiza procesul folosind GitHub Actions astfel încât workflowul
să aibă acces la clusterul local a fost nevoie de crearea unui GitHub runner local care să ruleze pe mașina de dezvoltare.

\subsection{Epuizarea memoriei GPU din cauza texturilor masive}
Un alt obstacol major a fost legat de modelele 3D ale Pământului și hărțile de fundal, care aveau rezoluții foarte mari (16K×8K).
Acestea produceau erori de memorie în WebGL (pierdere de context, erori la shadere și avertismente de redimensionare a texturilor). 
Problema a fost rezolvată prin utilizarea unor variante de 4K–8K pentru texturi, încărcare leneșă a acestora și folosirea formatelor
de texturi comprimate pentru a reduce consumul de memorie și a preveni pierderile de context WebGL.


\subsection{Crearea fluxului de parsare a codului C++}

Cea mai dificilă parte a fost design-ul antetului \texttt{tracer.h}: au trebuit definite
meta-trăsături (\textit{type traits}) care să distingă între \texttt{vector}, \texttt{list},
\texttt{map}, \texttt{priority_queue} sau chiar \texttt{std::array} și altele, fără ajutorul
reflecţiei native. În plus, a fost nevoie de serializarea stării fiecărui 
container. Asta a complicat puternic manipularea tipurilor compuse, precum 
\texttt{std::pair<K,V>} din map-uri. Implementarea funcţiei auxiliare
\lstinline{serializePair()} şi a ramurilor specializate din \lstinline{getState()}
/ \lstinline{getMetadata()} a solicitat multe iteraţii pentru a gestiona corect
conversiile către JSON. În final s-a putut evita ambiguitatea dintre clasele STL şi tipurile
primitive. Nu în ultimul rând, menţinerea consistenţei între operaţiile detectate
în \lstinline{TransformCode()} şi cele raportate de template-urile tracer a
cerut testare repetată, deoarece o singură pierdere (de exemplu
\lstinline{deque::shrink_to_fit}) genera stări incoerente.

\subsection{Dimensiunea clusterului Kubernetes}

Un alt obstacol a fost dimensiunea clusterului Kubernetes care este practic stocat 
prin imagini și volume Docker. Fără curățare constantă s-au ajuns la dimensiuni de aproape
60GB, ceea ce a dus la probleme de spațiu pe disc și la dificultăți în gestionarea
resurselor. Soluția a fost curățarea imaginilor nefolosite și compresarea volumelor Docker, iar
apoi compresarea întregii imagini de stocare wsl\footnote{https://learn.microsoft.com/en-us/windows/wsl/}
de mai multe ori pentru a putea continua dezvoltarea.

\subsection{Popularea și accesul extern la intrările in baza de date}
Pentru a crea entități noi sau a adauga date educaționale în MongoDB într-un pod
de Kubernetes, a fost necesară crearea unui job dedicat care să aibă acces în interiorul
clusterului, deoarece nu se poate accesa direct din afara clusterului cu cereri
directe. Asta a însemnat crearea unui fișier de configurare YAML pentru jobul respectiv și 
crearea unui pod dedicat pentru rularea acestuia.


\section{Analiza metricilor}
\label{subsec:analiza-metrici}

Cei mai reprezentativi indicatori de exploatare în contextul unei platforme ca VisioScience3D 
sunt timpii de răspuns, rata de cadre pe secundă (FPS) în scenele 3D, costurile de licențiere
și consumul de memorie RAM. Timpul de răspuns a fost constant sub 250ms pentru majoritatea cererilor
în mediul de dezvoltare, putând să difere în eventuala producție în funcție de încărcarea serverului.
Rata de cadre pe secundă este asigurată de tehnologiile WebGL și Three.js folosite și de 
tehnicile de optimizare incluse. În general se menține constantă la cel puțin 60 FPS.
Costurile de licențiere sunt zero pentru VisioScience3D, în timp ce alte platforme competitoare
implică taxe anuale semnificative. În ceea ce privește consumul de memorie RAM, platforma
se menține sub 150MB în majoritatea scenelor, ne încărcând resurse de foarte mari
dimensiuni.

Metricile au fost obținute prin utilizarea browserului Chrome pentru accesarea siteurilor și folosind informațiile
de latență la nivel de cerere, resursele utilizate raportate de browser și rata de cadre resimțită în scenele 3D
la utilizare (sub 60 de cadre pe secundă se resimt mici sacadări în animații și interacțiuni).

VisioScience3D răspunde în medie mai repede cel puțin în mediul de dezvoltare decât platforme gratuite 
şi depăşeşte nivelul minim de 60 cadre pe secundă, prag important pentru
experienţa vizuală fluidă. În testele de stres pe arhitectura backend-ului, aplicaţia a susţinut
până la sute de cereri concurente fără să rescalaze pod-urile folosite pentru servicii.
Din punct de vedere monetar, lipsa taxei de licenţiere o poziţionează peste Mozaweb şi IQBoard,
care introduc un cost de utilizare semnificativ.
Tehnicile adăugate ca optimizarea texturilor şi încărcarea leneșă ajută și mai mult la menţinerea
consumului de memorie sub 150 MB.

